<%# all TODOs should support following commands-%>
% TODO (and comment) configuration
%
% - \todo (from zebra, todonotes, luatodonotes), or \TODO - for "normal" TODOs
% - \todofix, or \FIXME - "important" TODOs
%
% - \textcomment - highlights text and has a hover comment
% - \sidecomment - just puts a comment to the side. Note: \comment is the comment environment of several packages (mathdesign, mindflow, verbatim, and others); only zebra adds \comment[name]{...} as command
%
% - \missingfigure
%
% - \textmarker
% - \modified
% - \change      - addresses a review comment

<% switch (todo) { case "zebra": -%>
% TODOs as short inline notes with a small numbered symbol in the margin; a summary and a list of all notes follow at the end
% Doc: https://ctan.org/pkg/zebra
% Commands: \todo[name]{...}, \note, \fixed, \placeholder, and \comment
\usepackage[draft]{zebra} % explicit: the class option draft=false would hide all notes
% Final version: hide all notes and the lists at the end
% \usepackage[final]{zebra}
% Important TODOs as own note type: \fixme{...}
\zebranewnote{fixme}{red!80!black}
\newcommand{\todofix}[1]{\zebrafixme{#1}}
\newcommand{\sidecomment}[1]{\zebracomment{#1}}
% \comment is also the comment environment (verbatim, comment; used by mindflow and acmart), which zebra does not redefine:
% \comment[name]{...} is a zebra note, \begin{comment} ... \end{comment} still comments out
\makeatletter
\AddToHook{begindocument/before}{%
  \@ifundefined{comment}{\let\comment\zebracomment}{%
    \let\ltg@commentenvironment\comment
    \def\ltg@commentname{comment}%
    \renewcommand*{\comment}{%
      \ifx\@currenvir\ltg@commentname
        \expandafter\ltg@commentenvironment
      \else
        \expandafter\zebracomment
      \fi}}}
\makeatother
\newcommand{\missingfigure}[1][]{\zebraplaceholder{missing figure}}

<% break; case "pdfcomment": -%>
% Enable nice comments
\usepackage{pdfcomment}

\ifdefined\directlua
  % Highlighting that breaks across lines (\colorbox does not): lua-ul's \highLight; the color needs luacolor
  \usepackage{luacolor,lua-ul}
  \newcommand{\textcomment}[2]{\highLight[yellow!60]{#1}\pdfcomment[color={0.234 0.867 0.211},hoffset=-6pt,voffset=10pt,opacity=0.5]{#2}}
\else
  \newcommand{\textcomment}[2]{\colorbox{yellow!60}{#1}\pdfcomment[color={0.234 0.867 0.211},hoffset=-6pt,voffset=10pt,opacity=0.5]{#2}}
\fi

% Small PDF comment
% 1st parameter: comment
\newcommand{\sidecomment}[1]{\pdfcomment[color={0.045 0.278 0.643},voffset=4pt,icon=Note]{#1}}
% Disabled variant - for the final PDF
%\newcommand{\sidecomment}[1]{}

\newcommand{\todo}[1]{TODO!\sidecomment{#1}}

% Changes
%
% 1st parameter: review comment
% 2nd parameter: new text
\newcommand{\change}[2]{{\color{red}#2}\pdfcomment[color={0.234 0.867 0.211},voffset=8pt,opacity=0.5]{#1}}
% Disabled variant - for the final PDF
%\newcommand{\change}[2]{#2}

<% break; case "todonotes": -%>
<% if (latexcompiler == "both") { -%>
\ifluatex
  % LuaLaTeX: luatodonotes places the notes so that notes and lines do not overlap
  \usepackage[color=gray!17]{luatodonotes}
\else
  \usepackage[color=gray!17]{todonotes}
\fi
<% } else if (latexcompiler == "lualatex") { -%>
% luatodonotes: todonotes with a placement that avoids overlapping notes and lines (LuaLaTeX only)
\usepackage[color=gray!17]{luatodonotes}
<% } else { -%>
\usepackage[color=gray!17]{todonotes}
<% } -%>
\let\xtodo\todo
<% if (isThesis) { -%>
% Inline TODOs as a titled box, which stands out in a long text
\usepackage{tcolorbox}
\newtcolorbox{todobox}{colback=yellow!8,colframe=brown!60!black,fonttitle=\bfseries,title=TODO}
\renewcommand{\todo}[1]{\begin{todobox}#1\end{todobox}}
<% } else { -%>
\renewcommand{\todo}[1]{\xtodo[inline,color=black!5]{#1}}
<% } -%>
\newcommand{\todofix}[1]{\xtodo[color=red!40]{#1}}

<% break; default: -%>
<%# everything defined below -%>

<% break; } -%>
% Define default commands
\makeatletter
\@ifundefined{missingfigure}{\newcommand{\missingfigure}{\ldots\ missing figure \ldots}}{}
\@ifundefined{textcomment}{\newcommand{\textcomment}[2]{#1 \todo{#2}}}{}
\@ifundefined{sidecomment}{\newcommand{\sidecomment}[1]{\marginpar{#1}}}{}
\@ifundefined{todo}{\newcommand{\todo}[1]{\sidecomment{#1}}}{}
\@ifundefined{TODO}{\newcommand{\TODO}[1]{\todo{#1}}}{}
\@ifundefined{todofix}{\newcommand{\todofix}[1]{\todo{#1}}}{}
\@ifundefined{FIXME}{\newcommand{\FIXME}[1]{\todofix{#1}}}{}
\@ifundefined{change}{\newcommand{\change}[2]{#1 $\rightarrow$ #2}}{}
\makeatother

% Textmarker (Textfarbe rot)
\newcommand{\textmarker}[1]{{\color{red} #1}\xspace}

% Modified (Text blau)
\newcommand{\modified}[1]{{\color{blue!60!black} #1}\xspace}

% Edits and comments by several authors, each with their own color and initials (phfcc)
% Doc: https://ctan.org/pkg/phfcc
\usepackage{phfcc}
\ifdefined\directlua
  \phfSetDefaultCommentingStyle{style=rmstrikethrough} % strike through removed text across lines (LuaLaTeX only)
\fi
% Comments as footnotes keep the layout intact, e.g., to check page limits
%\phfSetDefaultCommentingStyle{style=footcomments}
% One command per author: \FA{added text}, \FA*{text to remove}, \FA[a comment], \FA!{important};
% the initials appear once per page in the margin; the forms do not combine (\FA*[...]{...} is only a comment), so
% write \FA*{removed}\FA[comment]. Add one line per further author with their initials.
\phfMakeCommentingCommand[initials={FA}]{FA} % first author
% Final version: hide all edits and comments (there is no class option for it: "final" is the default of every class)
%\phfDisableCommentingCommands
